\chapter{Conclusions}
\label{ch:conclusions}

This thesis set out to answer one question: can the statistics of the one-dimensional Burridge--Knopoff model of an earthquake fault be
understood by mapping it onto a sandpile model, and can its exponents be obtained from the theory of self-organized criticality? This chapter
collects the answers, in the order of the aims listed in Chapter~\ref{ch:intro}, and lists the questions that remain open.

\section{Results}

\paragraph{Self-organized criticality in one dimension.} The Bak--Tang--Wiesenfeld sandpile is Abelian, which gives exact results in any
dimension: the mean number of topplings is the Green function of the toppling matrix, $\avg{s}\propto L^2$, and $D(2-\tau) = 2$. In one dimension
the model can be solved completely; its avalanche sizes are given by the closed formula $s = g(j)+g(l')-g(l)$, and their distribution is a function of
$s/L^2$ that is not a power law. The one-dimensional BTW model is therefore not critical, and whatever power laws the one-dimensional BK model produces
cannot come from the mechanism of the simplest sandpile (Chapter~\ref{ch:soc}).

\paragraph{The papers of Carlson and collaborators.} The 1989 and 1991 papers study different versions of the model: the 1989 version has a symmetric
friction law and a finite loading speed $\nu$, which sets the size of the smallest events; the 1991 version has an asymmetric law with a friction drop
$\sigma$ and $\nu\to0$. Their analytical estimates (the one-block event, slip pulses that grow exponentially at the sound speed, the localization length
$\tilde\xi$ and the crossover moment $\tilde M$) are the same in both, with $\nu$ replaced by $\sigma$. The one-block event, which underlies all of them,
is underestimated by a factor of about two for $\alpha=3$ when velocity weakening is neglected; we derived the corrected result,
$\delta U_1 = \sigma(1+e^{\pi\alpha/\omega})/\omega_1^2$ (Chapter~\ref{ch:bk}).

\paragraph{Implementation and validation.} The earlier implementation of this project contained several errors, among them a wrong Runge--Kutta stage, a coupling
constant tied to the system size, fixed instead of free ends, and event durations that were never recorded. The new programs (Python, with the time-critical loops
compiled by Numba) are validated by tests against exactly known results: the one-block event, a reference integrator, the energy balance, the sum rule, the Abelian
property and the exact one-dimensional solution of the BTW model (Chapter~\ref{ch:numerics}). With these programs, a run of the BK chain that records millions of events
takes minutes on a single processor core, so independent runs are executed in parallel rather than parallelising a single run.

\paragraph{Reproduction.} Both papers are reproduced. For the 1991 model the shape of the distributions and the exponents agree with the published values
($b = 0.91$--$0.96$ for $\alpha = 3$, against $0.92$; $b\approx0.36$--$0.40$ for $\alpha=1$, against $0.35$). The published curves lie, however, a factor of about five
below the normalisation defined in the paper, as the sum rule shows. For the 1989 model all five panels of Fig.~6 of Carlson and Langer are reproduced, with
$b = 0.89$--$0.94$ for $\alpha = 2.5$. Independent runs of the original length show that the shortest published runs could not determine $b$ better than about
$\pm0.1$; together with the event definition (space--time clusters of slip), this explains why the earlier version of this project failed to reproduce the 1989 figures
(Chapter~\ref{ch:results}).

\paragraph{The mapping.} Removing inertia from the BK equations gives a continuous sandpile on the chain with a fixed drop of force and bulk dissipation
$\epsilon = 1/(2\ell^2+1)$, or, with a different toppling rule, the one-dimensional Olami--Feder--Christensen model. The mapping is exact for the elementary process
(the one-block event) and for the force bookkeeping and the sum rule. It fails for everything else: the deterministic sandpile limit is exactly periodic or synchronised,
it is not chaotic, its avalanches are confined to a few stiffness lengths, and the friction parameters only set the unit of slip. Adding random thresholds removes the
periodicity but not the confinement. None of the features of the BK size distribution beyond the smallest events survives (Chapter~\ref{ch:mapping}).

\paragraph{Is the BK model critical?} Not in the strict sense of self-organized criticality. Its localized events follow a power law over a window bounded below by the
microscopic events and above by a crossover at a moment of a few units of $a$, and the great events form a separate population with a peak and a sharp cutoff. The friction
drop $\sigma$, which Carlson \emph{et al.} identify as the critical parameter, changes only the smallest events. The window of localized events and the range of great events
both widen as the stiffness ratio $\ell$ grows, while the exponent drifts slowly, so the continuum limit $\ell\to\infty$ remains an open possibility for scale invariance
(Chapter~\ref{ch:results}).

\paragraph{Approximate theory.} Mean-field theory (the branching process) describes the sandpile limit as a subcritical process with branching ratio $1-\epsilon$; it
predicts the mean avalanche size but not the distribution, which is instead described by a random walk with losses and collapses with the single scale $\ell$. The BK model
differs by inertial amplification: a rigid group of $n$ blocks slips by an amount that grows rapidly with $n$ and runs away above $n^* = 2\ell^2/(\alpha^2-1)$, and the
simulations show the same amplification. Great events require $\alpha>1$, the condition for a long group to run away. A sharper criterion, $n^*<\tilde\xi$, predicts their onset
correctly for $\ell = 6$ but was refuted by new runs at $\ell = 12$, where the onset moves to smaller $\alpha$ instead of larger (Chapter~\ref{ch:theory}). The predicted crossover moment $\tilde M = 2\xi/\alpha$ is not confirmed, and the exponent of the localized events depends on $\alpha$,
which rules out a universal mean-field exponent.

\section{The answer}

The one-dimensional Burridge--Knopoff model cannot be mapped onto a sandpile model. The obstacle is not a detail of the toppling rule or the absence of disorder, but the
inertia of the blocks: with velocity-weakening friction, a slipping region gains momentum and passes on more force the larger it becomes, so that its effective branching ratio
grows with its size, while in any sandpile built on static force transfer every toppling passes on the same fraction $1-\epsilon$ of a fixed drop. This amplification produces the
great events, the chaotic dynamics and the dependence on $\alpha$; its absence makes the sandpile limit subcritical and confined to the scale $\ell$. The Gutenberg--Richter-like
statistics of the BK model are therefore intrinsically dynamical, and they are not self-organized critical in the sense of the sandpile models.

\section{Open questions}

\begin{itemize}
  \item \textbf{The crossover.} The crossover from localized to great events lies at a moment of a few units of $a$, nearly independent of $\alpha$, $\sigma$ and $\ell$ in the range
    studied. A theory of its position, and a test whether it eventually scales with $\xi$ for $\ell\gg24$, are missing.
  \item \textbf{The runaway size.} The onset of great events lies at $\alpha\approx1.3$ for $\ell=6$ and below $1.25$ for $\ell = 12$. An estimate of the size at which an
    event runs away, based on the pulse solution rather than on a rigid group, should explain this dependence on $\ell$.
  \item \textbf{The exponent of the localized events.} It depends on $\alpha$, with a maximum $\tau\approx2$ near $\alpha = 3$. In the picture of Carlson \emph{et al.} it is set by
    the roughness that great events leave behind; a quantitative theory of that roughness would be needed to predict it.
  \item \textbf{The continuum limit.} Whether the window of the localized power law becomes infinite and its exponent converges as $\ell\to\infty$ requires chains with $\ell\gtrsim100$,
    which are within reach of the present programs with longer runs.
  \item \textbf{A dynamical cellular automaton.} A sandpile-like model that carries a velocity, or a rule in which the force passed on grows with the size of the slipping region,
    might reproduce the BK statistics at a fraction of the cost. The runaway of long slipping regions for $\alpha>1$, found in Chapter~\ref{ch:theory}, is the property such a model would have to reproduce.
  \item \textbf{Two dimensions.} The OFC model was derived from a two-dimensional BK model. Whether inertia plays the same role there, and whether the two-dimensional sandpile limit
    is closer to criticality, is the natural extension of this work.
\end{itemize}
